%%%%%%%%%%%%%%%%%%%%%%%%%%%%%%%%%%%%%%%%%%%%%%%%%%%
%% Exercices — Informatique Quantique — Pascal
%% Stage LIA / CERI — Université d'Avignon
%% Stagiaire : Patrice SEBASTIANO
%% Encadrant : Serigne GUEYE
%%%%%%%%%%%%%%%%%%%%%%%%%%%%%%%%%%%%%%%%%%%%%%%%%%%

% \documentclass[light]{ceri/sty/rapport}
% \documentclass[handout]{ceri/sty/rapport}
\documentclass{ceri/sty/rapport}


%%%%%%%%%%%%%%%%%%%%%%%%%%%%%%%%%%%%%%%%%%%%%%%%%%%
%% Informations générales
%%%%%%%%%%%%%%%%%%%%%%%%%%%%%%%%%%%%%%%%%%%%%%%%%%%
\major{Master d'Informatique}
\specialization{Informatique Quantique}
\course{Stage de Recherche — LIA Avignon}

\title{Recuit Adiabatique}
\subtitle{Hamiltonien de Rydberg \& Mapping QUBO --- Pasqal / Pulser}

\author{
	Patrice SEBASTIANO
}

\advisor[Encadrant]{
	Serigne GUEYE
}


%%%%%%%%%%%%%%%%%%%%%%%%%%%%%%%%%%%%%%%%%%%%%%%%%%%
%% Packages supplémentaires
%%%%%%%%%%%%%%%%%%%%%%%%%%%%%%%%%%%%%%%%%%%%%%%%%%%
\usepackage{ifthen}
\usepackage{amsmath, amssymb}
\usepackage{physics}
\usepackage{mathtools}
\usepackage{braket}
\usepackage{tcolorbox}
\tcbuselibrary{breakable, skins, listings}
\usepackage{listings}
\lstset{
  language=Python,
  basicstyle=\ttfamily\small,
  keywordstyle=\bfseries\color{blue!70!black},
  stringstyle=\color{green!50!black},
  commentstyle=\itshape\color{gray},
  breaklines=true,
  frame=none,
  showstringspaces=false,
}

% Couleurs personnalisées
\definecolor{exo-bg}{RGB}{230, 240, 255}
\definecolor{exo-border}{RGB}{50, 100, 200}
\definecolor{sol-bg}{RGB}{230, 255, 235}
\definecolor{sol-border}{RGB}{40, 160, 60}

% Environnement exercice — argument optionnel = référence page du livre
\newcounter{exercice}[section]
\renewcommand{\theexercice}{\thesection.\arabic{exercice}}

\newtcolorbox[use counter=exercice]{exercice}[1][]{%
	colback=exo-bg,
	colframe=exo-border,
	fonttitle=\bfseries,
	title={Exercice~\theexercice \ifthenelse{\equal{#1}{}}{}{~---~#1}},
	breakable,
	before upper={\parindent=1em}
}

\newtcolorbox{solution}{%
	colback=sol-bg,
	colframe=sol-border,
	fonttitle=\bfseries,
	title={Solution},
	breakable,
	before upper={\parindent=1em}
}


%%%%%%%%%%%%%%%%%%%%%%%%%%%%%%%%%%%%%%%%%%%%%%%%%%%
%% Début du document
%%%%%%%%%%%%%%%%%%%%%%%%%%%%%%%%%%%%%%%%%%%%%%%%%%%
\begin{document}

\maketitle
\chead{\textcolor{black}{Recuit Adiabatique -- Hamiltonien de Rydberg \& Mapping QUBO -- Pasqal}}
\sloppy

%%%%%%%%%%%%%%%%%%%%%%%%%%%%%%%%%%%%%%%%%%%%%%%%%%%
\section{Notes --- Hamiltonien adiabatique (Pulser / Pasqal)}
\label{sec:hamiltonien}
%%%%%%%%%%%%%%%%%%%%%%%%%%%%%%%%%%%%%%%%%%%%%%%%%%%

%-------------------------------------------------------
\subsection{L'Hamiltonien de Rydberg Général}
%-------------------------------------------------------

Sur une plateforme Pasqal, le registre d'atomes neutres est soumis à deux types d'interactions :
chaque atome est individuellement piloté par une impulsion laser, et les paires d'atomes excités
interagissent via les forces de Van der Waals. L'Hamiltonien total se décompose donc naturellement
en un terme de pilotage et un terme d'interaction :

\begin{equation}
H(t) = \sum_i H_i^D(t) \;+\; \sum_{j<i} H_{ij}^{\mathrm{int}}
\end{equation}

\paragraph{Terme de pilotage $H_i^D(t)$.}
Chaque atome est piloté par une impulsion laser décrite par trois paramètres : la fréquence de
Rabi $\Omega(t)$, le detuning $\delta(t)$, et la phase $\varphi$. L'Hamiltonien de pilotage
d'un atome unique sur ses deux niveaux $|g\rangle$ (fondamental) et $|r\rangle$ (Rydberg)
s'écrit, en unités de $\hbar$ (documentation Pulser) :

\begin{equation}
\frac{H^D(t)}{\hbar} = \frac{\Omega(t)}{2}\,e^{-i\varphi}\,|g\rangle\langle r|
\;+\; \frac{\Omega(t)}{2}\,e^{+i\varphi}\,|r\rangle\langle g|
\;-\; \delta(t)\,|r\rangle\langle r|
\end{equation}

En développant sur la base $\{|r\rangle, |g\rangle\}$, on reconnaît les matrices de Pauli :

\begin{equation}
\frac{\Omega(t)}{2}\,e^{-i\varphi}\,|g\rangle\langle r|
+ \frac{\Omega(t)}{2}\,e^{+i\varphi}\,|r\rangle\langle g|
= \frac{\Omega(t)}{2}
\begin{pmatrix} 0 & e^{i\varphi} \\ e^{-i\varphi} & 0 \end{pmatrix}
= \frac{\Omega(t)}{2}\big(\cos\varphi\,\sigma^x - \sin\varphi\,\sigma^y\big)
\end{equation}

Dans le protocole Pasqal pour le QUBO, le laser impose $\varphi = 0$ à tous les atomes,
ce qui annule le terme en $\sigma^y$. Le terme de pilotage se simplifie alors en
$\frac{\Omega(t)}{2}\sigma^x$.

\paragraph{Terme d'interaction $H_{ij}^{\mathrm{int}}$.}
Deux atomes simultanément dans l'état de Rydberg $|r\rangle$ ressentent une répulsion
dipôle-dipôle qui décroît comme la sixième puissance de leur distance $r_{ij}$ :

\begin{equation}
H_{ij}^{\mathrm{int}} = \frac{C_6}{r_{ij}^6}\,\hat{n}_i\hat{n}_j
\end{equation}

\paragraph{Hamiltonien explicite ($\varphi = 0$).}
En développant $H^D$ avec $\varphi = 0$ et en sommant sur tous les atomes, on obtient
l'expression complète (voir annexe~\ref{ann:formalisme} pour les détails sur la convention
des états et la définition matricielle de $\hat{n}_i$) :

\begin{equation}
H(t) = \underbrace{\Omega(t)\sum_i \sigma_i^x}_{\text{couplage Rabi}}
\;-\; \underbrace{\hbar \sum_i \delta_i(t)\,\hat{n}_i}_{\text{detuning local}}
\;+\; \sum_{j<i} \underbrace{\frac{C_6}{r_{ij}^6}}_{\text{interaction vdW}} \hat{n}_i \hat{n}_j
\end{equation}

%-------------------------------------------------------
\subsection{États Propres des Termes de l'Hamiltonien}
\label{sec:etats-propres}
%-------------------------------------------------------

Chaque terme de $H(t)$ admet un état propre de plus basse énergie. Identifier ces états
est la clé pour comprendre la dynamique du recuit adiabatique.

\medskip
\noindent\textbf{1. Terme de couplage Rabi : minimisé par la superposition uniforme.}

Le terme $\Omega(t)\sum_i \sigma_i^x$ agit via $\sigma^x$, dont les valeurs propres sont $\pm 1$ :

\begin{equation}
\sigma^x = \begin{pmatrix} 0 & 1 \\ 1 & 0 \end{pmatrix},
\qquad
\sigma^x\,|{+}\rangle = +|{+}\rangle,
\qquad
\sigma^x\,|{-}\rangle = -|{-}\rangle
\end{equation}

avec $|{\pm}\rangle = \frac{1}{\sqrt{2}}\big(|0\rangle \pm |1\rangle\big)$.
L'état propre de valeur propre $-1$ (le plus bas, puisque $\Omega > 0$) est $|{-}\rangle$.
Sur $n$ atomes, l'état qui minimise $\Omega\sum_i\sigma_i^x$ est le produit tensoriel
$|{-}\rangle^{\otimes n}$, qui n'est autre que la superposition uniforme de tous les états
de la base computationnelle :

\begin{equation}
|{-}\rangle^{\otimes n}
= \frac{1}{\sqrt{2^n}}\sum_{k=0}^{2^n-1}(-1)^{|k|}\,|k\rangle
\end{equation}

où $|k|$ désigne le poids de Hamming de $k$ (nombre de bits à $1$). En particulier, pour $n=2$ :

\begin{equation}
|{-}\rangle^{\otimes 2}
= \frac{1}{2}\big(|00\rangle - |01\rangle - |10\rangle + |11\rangle\big)
\end{equation}

\medskip
\noindent\textbf{2. Terme de detuning : minimisé par $|g\rangle^{\otimes n}$.}

Le terme $-\hbar\sum_i\delta_i(t)\,\hat{n}_i$ agit via le projecteur $\hat{n}_i = |r\rangle\langle r|_i$.
On vérifie d'abord sur un atome que $|g\rangle$ est vecteur propre de valeur propre $0$ :

\begin{equation}
\hat{n}\,|g\rangle
= \begin{pmatrix} 1 & 0 \\ 0 & 0 \end{pmatrix}
\begin{pmatrix} 0 \\ 1 \end{pmatrix}
= 0\,|g\rangle
\end{equation}

Sur $n$ atomes, l'opérateur $\hat{n}_i$ agit comme $I\otimes\cdots\otimes\hat{n}\otimes\cdots\otimes I$
(identité sur tous les sites sauf le site $i$). Son action sur $|gg\cdots g\rangle$ donne donc :

\begin{align}
\hat{n}_i\,|gg\cdots g\rangle
&= |g\rangle\otimes\cdots\otimes\underbrace{\hat{n}\,|g\rangle}_{=\,0}\otimes\cdots\otimes|g\rangle
= 0
\end{align}

La somme sur tous les sites s'annule terme à terme :

\begin{align}
-\hbar\sum_i\delta_i\,\hat{n}_i\;|gg\cdots g\rangle
= \sum_i\big(-\hbar\delta_i\big)\,\hat{n}_i\,|gg\cdots g\rangle
= 0
\end{align}

L'état $|gg\cdots g\rangle$ est donc vecteur propre de valeur propre nulle du terme de detuning.
Lorsque tous les coefficients $-\hbar\delta_i$ sont positifs (ce qui correspond à $\delta_i < 0$),
la valeur propre nulle est bien la plus basse possible : $|gg\cdots g\rangle$ est l'état qui
minimise ce terme. C'est précisément la condition au début du recuit adiabatique
(section~\ref{sec:recuit}, phase~1).

\medskip
\noindent\textbf{3. Terme d'interaction de Van der Waals : minimisé par $|g\rangle^{\otimes n}$.}

Le terme $\sum_{j<i}\frac{C_6}{r_{ij}^6}\,\hat{n}_i\hat{n}_j$ a des coefficients strictement
positifs ($C_6 > 0$). L'opérateur $\hat{n}_i\hat{n}_j$ agit comme
$I_1\otimes\cdots\otimes\hat{n}_i\otimes\cdots\otimes\hat{n}_j\otimes\cdots\otimes I_n$,
où les facteurs sont ordonnés selon les indices de site croissants $1 \le \cdots < i < \cdots < j < \cdots \le n$.
Son action sur $|gg\cdots g\rangle$ s'écrit donc, en respectant cet ordre :

\begin{equation}
\hat{n}_i\hat{n}_j\;|gg\cdots g\rangle
= |g\rangle_1 \otimes \cdots \otimes \underbrace{\hat{n}\,|g\rangle_i}_{=\,0}
  \otimes \cdots \otimes \underbrace{\hat{n}\,|g\rangle_j}_{=\,0}
  \otimes \cdots \otimes |g\rangle_n
= 0
\end{equation}

La somme sur toutes les paires $(i,j)$ s'annule donc terme à terme :

\begin{align}
\sum_{j<i}\frac{C_6}{r_{ij}^6}\,\hat{n}_i\hat{n}_j\;|gg\cdots g\rangle
&= \sum_{j<i}\frac{C_6}{r_{ij}^6} \cdot 0 \;=\; 0
\end{align}

Puisque tous les coefficients $C_6/r_{ij}^6$ sont strictement positifs, la valeur propre nulle
est bien la plus basse atteignable : $|gg\cdots g\rangle$ est l'état fondamental du terme
d'interaction de Van der Waals. Combiné au résultat précédent sur le terme de detuning,
on conclut que $|gg\cdots g\rangle$ annule simultanément les deux derniers termes de $H(t)$,
ce qui en fait l'état de départ naturel du recuit adiabatique (section~\ref{sec:recuit}, phase~1).

%-------------------------------------------------------
\subsection{Mapping avec le Problème QUBO}
%-------------------------------------------------------

Un problème QUBO (Quadratic Unconstrained Binary Optimization) s'écrit sous la forme standard :

\begin{equation}
\min_{x} \;\sum_{i=1}^{N} c_i\, x_i \;+\; \sum_{i=1}^{N}\sum_{j=i+1}^{N} q_{ij}\, x_i x_j,
\qquad x_i \in \{0,1\},\quad i = 1,2,\dots,N
\end{equation}

Le mapping physique repose sur la correspondance terme à terme avec l'Hamiltonien de Rydberg.

\paragraph{Termes linéaires $\leftrightarrow$ detuning local.}
Puisque $\hat{n}_i$ joue le rôle de la variable binaire $x_i$, on identifie le coefficient
linéaire $c_i$ au detuning local $\delta_i$ via la relation :

\begin{equation}
-\hbar\,\delta_i(t) = c_i
\end{equation}

\paragraph{Termes quadratiques $\leftrightarrow$ interaction de Van der Waals.}
Le couplage $q_{ij}$ entre les variables $x_i$ et $x_j$ est encodé par l'interaction dipôle-dipôle
entre deux atomes de Rydberg séparés par une distance $r_{ij}$ :

\begin{equation}
q_{ij} \;=\; \frac{C_6}{(r_{ij})^6}
\end{equation}

En ajustant la géométrie du registre d'atomes (les positions $r_{ij}$), on contrôle directement
les poids quadratiques du QUBO. L'Hamiltonien de Rydberg est ainsi un encodage physique naturel
de la matrice QUBO.

%-------------------------------------------------------
\subsection{Protocole du Recuit Adiabatique}
\label{sec:recuit}
%-------------------------------------------------------

Le recuit adiabatique repose sur le \textbf{théorème adiabatique} : si un système quantique
est initialement dans l'état fondamental de son Hamiltonien, et si ce Hamiltonien évolue
suffisamment lentement, alors le système reste à tout instant dans l'état fondamental instantané.
Autrement dit, le système suit continûment l'état qui minimise l'énergie, sans jamais « décrocher »
vers un état excité.

C'est cette propriété qui rend le recuit adiabatique utile pour l'optimisation : en construisant
un Hamiltonien $H(t)$ dont l'état fondamental à $t = T$ encode la solution du QUBO, et en faisant
évoluer le système assez lentement depuis un état fondamental connu à $t = 0$, on garantit qu'à
la fin du protocole les atomes se trouvent dans l'état qui minimise l'énergie du problème. Il
suffit alors de les mesurer pour lire la solution optimale.

Le recuit se déroule en trois phases successives, dictées par l'évolution conjointe de
$\Omega(t)$ et $\delta_i(t)$ sur l'intervalle $[0, T]$.

\paragraph{Phase 1 --- État initial ($t = 0$, $\Omega = 0$, $\delta_i < 0$).}

Au début du protocole, le couplage Rabi est nul et le detuning est strictement négatif pour tout
site~$i$. L'Hamiltonien se réduit alors à ses deux derniers termes :

\begin{equation}
H(0) = -\hbar \sum_i \delta_i\,\hat{n}_i \;+\; \sum_{j<i} \frac{C_6}{r_{ij}^6}\,\hat{n}_i \hat{n}_j
\end{equation}

Les deux contributions sont de signe défini positif : le terme d'interaction Van der Waals a des
coefficients $C_6/r_{ij}^6 > 0$, et le terme de detuning vaut $-\hbar\delta_i > 0$ puisque
$\delta_i < 0$. L'énergie totale est donc une somme de termes strictement positifs pondérés par
les $\hat{n}_i$. Elle est minimisée de manière triviale en annulant toutes les variables
binaires, c'est-à-dire en plaçant tous les atomes dans l'état fondamental :

\begin{equation}
|\psi(0)\rangle = |gg\cdots g\rangle
\end{equation}

Le système est en équilibre adiabatique dans cet état, ce que confirment les calculs établis
à la section~\ref{sec:etats-propres} : les termes de detuning et d'interaction de Van der Waals sont
simultanément annulés par $|gg\cdots g\rangle$.

\paragraph{Phase 2 --- Exploration quantique ($\Omega = \Omega_{\max}$, $\delta_i \approx 0$).}

Au cours du recuit, $\Omega(t)$ monte jusqu'à son maximum tandis que $\delta_i(t)$ passe par zéro.
Le terme de couplage Rabi domine alors. Comme établi à la section~\ref{sec:etats-propres}, ce terme
est minimisé par $|{-}\rangle^{\otimes n}$, état propre de $\sigma^x$ de valeur propre $-1$ sur chaque
site, qui correspond à une superposition cohérente de toutes les configurations du registre. Cependant, le terme d'interaction de
Van der Waals reste présent ($C_6/r_{ij}^6$ ne dépend pas de $\delta$) et vient perturber cette
superposition idéale, créant de l'intrication entre les atomes. C'est cette combinaison —
superposition cohérente induite par $\Omega$ et intrication induite par les interactions de Van der
Waals — qui permet au registre d'explorer simultanément toutes les configurations possibles du
système. Pour une analyse détaillée de cet état intermédiaire, on renvoie le lecteur à la
littérature sur les protocoles QAA (Quantum Adiabatic Algorithm).

\paragraph{Phase 3 --- Lecture de la solution ($t = T$, $\Omega = 0$).}

En fin de protocole, $\Omega$ redescend à zéro tandis que les $\delta_i(T)$ ont atteint les
valeurs cibles fixées par le problème QUBO, via la correspondance $-\hbar\,\delta_i(T) = c_i$.
Les coefficients $c_i$ peuvent être positifs ou négatifs selon les données du problème.
L'Hamiltonien se réduit à ses deux derniers termes :

\begin{equation}
H(T) = -\hbar \sum_i \delta_i\,\hat{n}_i \;+\; \sum_{j<i} \frac{C_6}{r_{ij}^6}\,\hat{n}_i \hat{n}_j
\end{equation}

En substituant $-\hbar\delta_i = c_i$ et $C_6/r_{ij}^6 = q_{ij}$, on reconnaît terme à terme
la fonction objectif du QUBO, avec $\hat{n}_i$ jouant le rôle de la variable binaire $x_i$ :

\begin{equation}
H(T) \;\longleftrightarrow\; \sum_i c_i\,x_i \;+\; \sum_{j<i} q_{ij}\,x_i x_j
\end{equation}

Par le théorème adiabatique, le système se trouve à $t = T$ dans l'état fondamental de $H(T)$,
c'est-à-dire l'état qui minimise cette énergie. La mesure du registre donne alors directement
le bitstring $x^* \in \{0,1\}^n$ qui minimise le QUBO.



%%%%%%%%%%%%%%%%%%%%%%%%%%%%%%%%%%%%%%%%%%%%%%%%%%%
%% Annexes
%%%%%%%%%%%%%%%%%%%%%%%%%%%%%%%%%%%%%%%%%%%%%%%%%%%
\appendix

\section{Formalisme des États et du Projecteur $\hat{n}_i$}
\label{ann:formalisme}

\subsubsection*{Convention des états (Pulser)}

Dans la base Pulser, l'état fondamental est étiqueté $|1\rangle$ et l'état de Rydberg excité est
étiqueté $|0\rangle$ :

\begin{align}
|g\rangle &= |1\rangle = \begin{pmatrix}0\\1\end{pmatrix} \\[4pt]
|r\rangle &= |0\rangle = \begin{pmatrix}1\\0\end{pmatrix}
\end{align}

\subsubsection*{Définition du projecteur $\hat{n}_i$ et lien avec $\sigma^z$}

L'opérateur de comptage $\hat{n}_i$ est le projecteur sur l'état de Rydberg de l'atome $i$ :

\begin{equation}
\hat{n}_i \;=\; |r\rangle\langle r|_i
\;=\; \begin{pmatrix} 1 \\ 0 \end{pmatrix}\begin{pmatrix} 1 & 0 \end{pmatrix}
\;=\; \begin{pmatrix} 1 & 0 \\ 0 & 0 \end{pmatrix}
\end{equation}

Il s'exprime également à partir de la matrice de Pauli $\sigma^z$ :

\begin{equation}
\hat{n}_i = \frac{1+\sigma_i^z}{2}, \qquad
\sigma^z = \begin{pmatrix} 1 & 0 \\ 0 & -1\end{pmatrix}
\end{equation}

On vérifie cette identité matriciellement :

\begin{equation}
\frac{1+\sigma_i^z}{2} = \frac{1}{2}\left[
\begin{pmatrix}1&0\\0&1\end{pmatrix} + \begin{pmatrix}1&0\\0&-1\end{pmatrix}
\right]
= \frac{1}{2}\begin{pmatrix}2 & 0\\0&0\end{pmatrix}
= \begin{pmatrix}1&0\\0&0\end{pmatrix}
\end{equation}

\subsubsection*{Action de $\hat{n}$ sur les états de base}

L'opérateur $\hat{n}$ admet deux valeurs propres : $1$ pour l'état excité $|r\rangle = |0\rangle$
et $0$ pour l'état fondamental $|g\rangle = |1\rangle$ :

\begin{align}
\hat{n}\,|0\rangle &= 1\times|0\rangle \\
\hat{n}\,|1\rangle &= 0\times|1\rangle
\end{align}

\subsubsection*{Action de $\hat{n}_i\hat{n}_j$ sur une superposition générale}

Pour une superposition uniforme $|\psi\rangle = \frac{1}{\sqrt{2^n}}\sum_{k=0}^{2^n-1}|k\rangle$,
l'opérateur $\hat{n}_i\hat{n}_j$ agit terme à terme sur chaque état de base, les facteurs étant
ordonnés selon les indices de site croissants :

\begin{equation}
\hat{n}_i \hat{n}_j\,|\psi\rangle
= \frac{1}{\sqrt{2^n}}\sum_{k=0}^{2^n-1} \left(\hat{n}_i \otimes \hat{n}_j\right)|k\rangle
\end{equation}

Seuls les états $|k\rangle$ dont les bits $i$ et $j$ valent simultanément $1$ (i.e.\ les deux
atomes sont dans $|r\rangle$) contribuent de façon non nulle à cette somme.


%%%%%%%%%%%%%%%%%%%%%%%%%%%%%%%%%%%%%%%%%%%%%%%%%%%
%% Fin du document
%%%%%%%%%%%%%%%%%%%%%%%%%%%%%%%%%%%%%%%%%%%%%%%%%%%
\end{document}
